\documentclass[11pt]{article}
\usepackage[margin=1in]{geometry}
\usepackage{parskip}
\usepackage{xcolor}
\newcommand{\status}[1]{\par\noindent Status: #1\par\medskip\noindent{\color{gray!40}\hrulefill}\par\medskip}

\begin{document}


\section{Data and descriptive}

If worthwhile for some of these descriptive exercises, I have a PhD student who might be able to do them. We would have to get her access to the VRDC...

\status{No need to get another person involved. Plus it would cost \$15k for her to have a seat! Most of this is easy enough to code up so I'm working on it..}

Heterogeneity within DRG: distribution of patient risk scores.

\status{Pending. Distribution of the LPM risk score, for DRG 470 and the other DRGs.}

How well do the patient risk scores predict failure, net of the doc's observed performance?

\status{Pending. LPM of failure on the risk score with and without surgeon fixed effects.}

(R2 \#3) Volume - outcome relationship. Does performance increase with volume (over long-term) within surgeons? Does it decrease (over short-term, eg within a year) when there is high congestion? Could estimate LPMs with patient risk-adjusters and surgeon FEs.

\status{Pending. Within-surgeon LPM of failure on cumulative and current-quarter volume.}

CF reallocation and volume. Plot histogram of observed patient counts (e.g., per year), and plot in CF reallocation. Check the upper tails.

\status{Pending. Baseline against counterfactual volume per surgeon, with the upper tail.}

Show that surgeon volume (and its heterogeneity across docs) is pretty consistent over time. Could compute within vs between variance in volume (e.g., by month or quarter) and compare the two. (From R2's 3b criticism of our long time periods for congestion.)

\status{Pending. Within versus between variance of surgeon volume, quarterly and yearly.}

Heterogeneity in patient char's across surgeons. Plot dist's across surgeons of key demographics: race, sex, age (e.g., use averages of these char's over the 6 year estimation period). Also dist of risk scores. Then est correlation of these features with the observed quality (i.e., performance over 11 years).

\status{Pending. Patient mix by surgeon, observed against random reassignment, and its correlation with quality.}

Could we look at hospital discharge (or all payer) data in some state(s), to see how much volume we are capturing, and how much heterogeneity there is across surgeons? Concern is that best docs are at capacity (more than our descriptive analysis suggests) due to commercial patients. (R4 \#5)

\status{Not started; we could do this with state inpatient data, but not all states make this easy to get. Florida is pretty easy although it would cost a bit of money (not too much...likely in the hundreds). There's also aggregate stats out there showing the Medicare accounts for well over half of all joint replacements, but I'm not sure if those stats are traditional Medicare or add Medicare Advantage.}

\section{More programming and computation}

Other change to event study: DDD-ish comparing change in referrals (per fixed calendar time like a qtr) when see a failure vs when see a success. This is R2's idea to address mean reversion.

\status{Pending. Calendar-time event study, failures against successes.}

Could a surgeon-level event study include spillover from surgeon A to surgeon B, if A is in B's group? (R2 \#5)

\status{Pending. Same-group spillover as a second outcome.}

Should we also do a practice-level event study?

\status{Pending. Re-running the reduced form at the practice level.}

Logit interacting patient risk scores with surgeon average quality (maybe leave-one-out or 5 year avg before patient t, not pair-specific but all patients). R2 suggested to interact with distance, but this interaction with observed performance seems better.

\status{Pending. Logit with the risk score interacted with running surgeon quality.}

Logit using PCP movers?

\status{Pending. Pooled logit on movers.}

Logit using relatively new PCPs?

\status{Pending. Pooled logit on new PCPs.}

(above two could pool HRRs to increase power)

\status{Pending. Both pooled.}

Simulate the model for one market, using our estimated parameters. Program the full solution as well as Gittins. See how much the simulated behavior differs btw the two methods.

\status{Built locally in R (exact solution of the learning problem against the Gittins index). Hoping to run in the next few days.}

Not now - maybe, depending on some descriptive assessments:

(Estimate structural model using practice groups. R4 \#2and6)

\status{Pending. Full parallel set at the practice level. The earlier sign flip on the existing-familiarity counterfactual still needs explaining, and the new descriptives are aimed at that.}

(Include average effect of volume / surgeon learning in the structural model.)

\status{Not started. The within-surgeon LPM above is the descriptive check that would come first.}

\section{Theory}

\section{Other items from the reports}

Items we have taken up that are not on the list above.

R1 \#4--5, R2 \#2, R3 2a (the event study). \status{Done. Replaced by the re-referral design, now Section 2 of the paper (Table 3 and Figure 4), with the window construction and the worked example in the appendix.}

R3 1a (the Gittins index with logit shocks). \status{Pending. The simulation above is the test of this.}

R2 3a (the congestion term). \status{Writing. We can think of congestion as wait time (which we can't observe but can proxy for with a surgeon's total volume), so congestion enters as a disutility the same way distance does. If we mention this explicitly in the paper, then it might deal with some general concerns that congestion is kind of a reduced-form element introduced into a more structural model.}

R2 3c, R3 1b (the counterfactual fixed point). \status{Done. New closing paragraph in the Estimation Procedure subsection (5.3): the counterfactual iterates on surgeon volumes from the observed values until volumes and referral probabilities are consistent, with a pointer to the appendix check from two starting values.}

R2 3b (congestion at three-year windows). \status{Pending: I'm re-running with yearly congestion but only for some markets where specialists have enough referrals per year. We'll see what that shows. Otherwise we already show that the FEs don't move much from one 3yr period to another, so that could be something to add.}

R2 minor 2 (established PCPs against all PCPs). \status{Pending. Considering results for established versus all PCPs.}

R2 minor 3 (convergence). \status{Pending. Convergence now logged for every market so we can see what's happening.}

R2 minor 4 (myopic against forward-looking). \status{Pending. Vuong test by market to assess improvement in likelihood between the two.}

\section{Reviewer 1}

1. The identification claim against Ching, Erdem, and Keane (2013). \status{Done. The identification paragraph in the introduction is cut. The Identification subsection (4.4) now opens as a discussion of the myopic model's parameters with a conjecture for the forward-looking model, and its closing paragraph describes the sources of variation along the lines of Ching et al.\ with the formal-versus-verbal contrast dropped.}

2. DRG 470 and case mix. \status{Already in the appendix (D.4, D.5)...not sure there's anything else to do here.}

And all the work above regarding risk adjustment and heterogeneity across surgeons in case mix. Should we consider estimating the stuctural model with only DRG 470? We can show plenty of heterogeneity in performance within this one DRG.

3. Risk aversion as an alternative reading of the familiarity effect. \status{Done. Footnote at the definition of familiarity in Model Specification (4.1): with a binary outcome, expected utility is linear in the success probability, so there is no curvature through which risk aversion could enter; the familiarity effect is a direct preference for surgeons the PCP has used before.}

4. The gap between the event study (16\%) and the structural partial effect, and the suggestion to estimate the model on the event-study sample. \status{Already addressed in new reduced form analysis.}

I'm interested in the idea that we should use the same sample for DiD and structural model.

5. One failure should move a well-established posterior little; adaptive or recency-weighted learning. \status{Partly addressed, partly writing. Not sure about "adaptive" or "recency-weighted"...}

Is there a way to cleanly distinguish between the two? Can't use a single-patient shock as in the event study.

Is the adaptive variable entirely observable? If so, could we correlate it with the estimated Bayesian posterior mean?

6. Whether surgeons' skills evolve over the panel. \status{Incorporated as part of new analysis looking at within-surgeon probability of failure as a function of cumulative experience.}

If we are checking this, do we also want to check whether familiarity affects outcomes?

7. Positioning: intermediary learning is not new; patient-level outcomes together with the intermediary's choice is. \status{Done. The contribution paragraph in the introduction now says that physicians learning on behalf of patients is the setting of the prescribing literature, and that what is distinctive here is observing the outcome of each referred patient together with the PCP's subsequent choices.}

I don't know if the intermediary thing helps us.

\end{document}
